% =================================================
% MODELO DE RESUMO EXPANDIDO - PEPICT / UNIFESP
% =================================================

\documentclass[12pt,a4paper]{article}

% -------------------------------------------------
% PACOTES
% -------------------------------------------------

\usepackage[utf8]{inputenc}
\usepackage[T1]{fontenc}
\usepackage[brazil]{babel}
\usepackage{lmodern}
\usepackage{geometry}
\usepackage{setspace}
\usepackage{titlesec}

% -------------------------------------------------
% CONFIGURAÇÕES ABNT
% -------------------------------------------------

\geometry{
    left=3cm,
    right=2cm,
    top=3cm,
    bottom=2cm
}

\setlength{\parindent}{1.25cm}

\onehalfspacing

% -------------------------------------------------
% INÍCIO
% -------------------------------------------------

\begin{document}

% -------------------------------------------------
% CABEÇALHO INSTITUCIONAL
% -------------------------------------------------

\begin{center}

{\large \textbf{UNIVERSIDADE FEDERAL DE SÃO PAULO -- UNIFESP}}

\vspace{0.3cm}

{\normalsize PROGRAMA PEPICT}

\vspace{0.3cm}

{\normalsize ICT -- Instituto de Ciência e Tecnologia}

\vspace{1cm}

{\Large \textbf{titulo do trabalho}}

\end{center}

\vspace{1cm}

% -------------------------------------------------
% AUTORES
% -------------------------------------------------

\noindent
\textbf{Autores(as)}

Aluno Teste2

\vspace{0.7cm}

% -------------------------------------------------
% DISCIPLINA E CURSO
% -------------------------------------------------

\noindent
\textbf{Disciplina e Curso}

Arquitetura e Organização de Computadores

\vspace{0.7cm}

% -------------------------------------------------
% PROGRAMA PEPICT
% -------------------------------------------------

\noindent
\textbf{Programa PEPICT vinculado}

Educação e Sustentabilidade, Tecnologia e Inovação para o Desenvolvimento Social

\vspace{0.7cm}

% -------------------------------------------------
% OBJETIVOS
% -------------------------------------------------

\section*{Objetivos do Trabalho}

sistema deve permitir que o coordenador filtre turmas por semestre, visualize projetos aprovados, selecione os trabalhos que farão parte da revista e gere o arquivo PDF consolidado da edição.

% -------------------------------------------------
% METODOLOGIA
% -------------------------------------------------

\section*{Metodologia}

Concluída a etapa de requisitos, o passo seguinte foi estruturar o modelo de resumo. Para isso, realizamos reuniões com a orientação do projeto para alinhar e selecionar quais informações seriam essenciais e pertinentes para compor esses resumos estruturados, o que resultou no modelo ilustrado na imagem abaixo:

% -------------------------------------------------
% RESULTADOS
% -------------------------------------------------

\section*{Resultados e Produtos}

Foi desenvolvida em 2001, quando seus autores buscaram solucionar lacunas e ineficiências do desenvolvimento em cascata, detectadas quando essa metodologia não produzia mais resultados satisfatórios frente a vários parâmetros de medida. Era comum, por exemplo,  demorar anos entre a identificação de uma necessidade da empresa e a entrega do sistema pronto. Nesse meio-tempo, as mudanças no mercado e nas exigências do negócio faziam com que grande parte do projeto fosse cancelada antes mesmo de ser entregue. Esse desperdício de tempo e dinheiro motivou vários desenvolvedores a buscarem uma alternativa.

% -------------------------------------------------
% ODS
% -------------------------------------------------

\section*{Relação com os ODS}

ODS 2 - Fome Zero e Agricultura Sustentável, ODS 8 - Trabalho Decente e Crescimento Econômico, ODS 13 - Ação Contra a Mudança Global do Clima

% -------------------------------------------------
% REFLEXÃO CRÍTICA
% -------------------------------------------------

\section*{Reflexão Crítica }

A metodologia ágil é uma abordagem de desenvolvimento que visa criar um fluxo de trabalho para tornar o processo de produção o mais eficiente e assertivo possível. Nessa abordagem, o desenvolvimento é um processo colaborativo no qual todas as partes mantêm um diálogo claro, a fim de manter o produto o mais próximo do desejado pelo solicitante, dentro da realidade de desenvolvimento e recursos disponíveis

% -------------------------------------------------
% REFERÊNCIAS
% -------------------------------------------------

\section*{Referências}

A metodologia ágil é uma abordagem de desenvolvimento que visa criar um fluxo de trabalho para tornar o processo de produção o mais eficiente e assertivo possível. Nessa abordagem, o desenvolvimento é um processo colaborativo no qual todas as partes mantêm um diálogo claro, a fim de manter o produto o mais próximo do desejado pelo solicitante, dentro da realidade de desenvolvimento e recursos disponíveis

\end{document}